\scanpage{20}
\mspage{5)}
\uline{Proposition} Considérons la cat. des
systèmes $(X,Y,(X_i)_{i\in I}\deleted{\unsure{\}}},T)$ où
$X$ espace top., $Y,(X_i)_{i\in I}$ des sous-espaces fermés,
$T$ voisinage cônique de $\unsure{Y|X_i}$
\deleted{compatible}\inserted{compatible} avec
$(X_i)$. (pour les iso), et la catégorie
des systèmes
\[
(X^*,\dot X,(X_i^*),\deleted{\unsure{\}}}Y,
\deleted{\unsure{f}}\dot X\unsure{\longrightarrow|\mathrel{\relbar\joinrel\relbar}}Y)
\]
$X^*,Y$ espaces top\textsuperscript{s}, $\dot X$ et $X_i^*$, parties fermées de $X$,
$\dot X\longrightarrow Y$ application continue. Ces deux
catégories sont équivalentes.
\inserted{Si $T$ voisinage cônique strict}
Pour que $T$
soit \uline{strictement} \deleted{compatible avec les}
\inserted{transversal aux} $X_i$,
il f. \unsure{et|t} s. que $\dot X$ soit un sous-espace bord
de $\unsure{\dot X^*|X^*}$, \deleted{\uline{\unsure{strict}}} compatible avec les $X_i^*$.
Pour que $T$ soit équisingulier pour
$(X_i)$, il f. \unsure{et|t} s. que
\diagram{\[
\begin{tikzcd}
\dot X\arrow[r]&Y,\\
(\dot X_i)\arrow[u,phantom,"\cup" description]&
\end{tikzcd}
\]}
soit une fibration localem\textsuperscript{t} triviale
simultanée (où on pose $\dot X_i=X_i\cap\dot X$)
\margin{19}
